\subsection{The full two-adic Selmer image of the Fermat-eight Jacobian}

\begin{statement}{A.187}{Proved}
{Exact elliptic isogeny Selmer images, finite Prym Tate--Shafarevich
groups, and the global two-adic Kummer image}
For each of the three rational elliptics and the three
quadratic-field elliptics in A.185, both displayed dual
two-isogeny Selmer groups equal their actual connecting
images on rational points over the stated field.  In particular
\[
                         \Sha(E_b/k)[2^\infty]=0
\tag{A.187.1}
\]
for each of these six curves, and for its dual isogeny curve.
For the three rational Prym surfaces of A.182,
\[
  \Sha(\mathcal P_{A,B}/\Q)[2^\infty]
       =\Sha(\mathcal P_{A,B}/\Q)[2]
       \quad\text{is finite}.
\tag{A.187.2}
\]
Let \(F:J_8\to\mathcal B\) be the particular isogeny
(A.182.3), with exactly its displayed divisor-pushforward
components, and put
\(N=\deg F\), \(a=v_2(N)+1\).  Then
\[
 2^a\Sha(J_8/\Q)[2^\infty]=0,
 \qquad\Sha(J_8/\Q)[2^\infty]\text{ is finite},
 \qquad T_2\Sha(J_8/\Q)=0.
\tag{A.187.3}
\]
This uses the degree of the specified isogeny; no numerical
value of that degree is asserted.
For every \(n\geq1\), the exact finite comparison is
\[
 \operatorname{im}\bigl(
    \operatorname{Sel}^{2^{n+a}}(J_8/\Q)
         \longrightarrow\operatorname{Sel}^{2^n}(J_8/\Q)
                       \bigr)
 =\kappa_{2^n}\bigl(J_8(\Q)/2^nJ_8(\Q)\bigr).
\tag{A.187.4}
\]
The transition is the one induced by multiplication by
\(2^a\) on finite torsion coefficients.  Taking the actual
continuous cohomology comparison gives an isomorphism
\[
 \widehat{J_8(\Q)}_2
       \xrightarrow{\ \kappa_2\ }
 \operatorname{Sel}(\Q,T_2J_8)
 :=\{\eta\in H^1_{\rm cts}(\Q,T_2J_8):
       \eta\bmod2^n\in\operatorname{Sel}^{2^n}
                         \text{ for all }n\}.
\tag{A.187.5}
\]
Its source has free \(\Z_2\)-rank fourteen.  These maps
commute with the isogeny components in A.182 and every
localization map.

Consequently, the compact Jacobian two-adic image of a
full Selmer section of a nonempty open of \(C_8\), with
origin \(P_0=(1,1)\), is the Kummer class of one element
of \(\widehat{J_8(\Q)}_2\).  This specifies its global
two-adic image; it does not classify those sections or
identify an ordinary local point with a rational point.
\end{statement}

\begin{proof}
\emph{The degree-two torsors and the isogeny orientations.}
Retain the actual maps
\(\phi_b:E_b\to E_{-4b}=E_b'\) and
\(\psi_b:E_b'\to E_b\) in (A.185.8).
For a squareclass \(d\in k^\times/k^{\times2}\), the
proper quartic
\[
 C_d:\quad N^2=dU^4+(b/d)V^4,
 \qquad
 h_d:C_d\longrightarrow E_b,
 \quad (U,V,N)\longmapsto
       \left(d\frac{U^2}{V^2},d\frac{UN}{V^3}\right)
\tag{A.187.6}
\]
is the actual twist of the \(\psi_b\)-cover with class
\([d]\in H^1(k,E_b'[\psi_b])\simeq k^\times/k^{\times2}\).
To verify this assertion including its orientation, write
\(t=U/V,n=N/V^2,s^2=d\).  Over \(k(s)\) the map
\[
 C_d\longrightarrow E_b',\qquad
 X=2dt^2+2sn,\qquad Y=2stX
\tag{A.187.7}
\]
is an isomorphism with inverse
\(t=Y/(2sX),\ n=(X-2dt^2)/(2s)\).
Its composition with \(\psi_b\) is \(h_d\).
Replacing \(s\) by \(-s\) changes \(X\) to
\(-4b/X\) and changes \(Y\) to \(4bY/X^2\),
which is exactly translation by the kernel point
\(T'=(0,0)\) on \(E_b'\).  The formulas extend to
the proper curves and prove the claimed torsor and cocycle
comparison.  In particular \(C_d(k_v)\ne\varnothing\)
is precisely the local connecting-image condition for
the \(\psi_b\)-Selmer group.

Every upper bound in A.185 was obtained from such local
conditions: the valuation parity outside the coefficient
support, the real signs, the primitive quartic tests at
three, and the stated dyadic tests for the rational curves.
The endpoints \(U=0\) or \(V=0\) are included in the
primitive tests, so no torsor point is lost.  The localized
ideal-class arguments there applied to every squareclass
with those valuation parities, not only to images of
global points.  Thus the same lists bound the full
isogeny Selmer groups.  Since displayed global points
attain every surviving class, the two exact equalities are
\[
\begin{aligned}
 \operatorname{Sel}^{\psi_b}(E_b'/k)
     &=\alpha_b(E_b(k)),\\
 \operatorname{Sel}^{\phi_b}(E_b/k)
     &=\alpha_{-4b}(E_b'(k)).
\end{aligned}
\tag{A.187.8}
\]
The scalings of the dual curves in A.185 multiply the
first coordinate by a square, so they preserve these
literal squareclass maps.

For reference, the first and second groups in (A.187.8)
are generated as follows, where \(s\) has the row-specific
square and every generator denotes its squareclass:
\[
\begin{array}{c|c|c|c}
k&b&\operatorname{Sel}^{\psi_b}(E_b'/k)
       &\operatorname{Sel}^{\phi_b}(E_b/k)\\ \hline
\Q&-6&\langle3,-6\rangle&\langle6\rangle\\
\Q&18&\langle6,2\rangle&\langle-2\rangle\\
\Q&-12&\langle6,-3\rangle&\langle3\rangle\\
\Q(s),\ s^2=-3&-4s&\langle-2,-s\rangle&\langle s\rangle\\
\Q(s),\ s^2=2&6s&\langle3,3s\rangle&\langle3,-3s\rangle\\
\Q(s),\ s^2=-6&8s&\langle-2,2s\rangle&\langle-2,-2s\rangle.
\end{array}
\tag{A.187.9}
\]

The two finite isogeny Kummer sequences have the exact
orientations
\[
\begin{aligned}
0\longrightarrow E_b(k)/\psi_bE_b'(k)
 &\longrightarrow\operatorname{Sel}^{\psi_b}(E_b'/k)
 \longrightarrow\Sha(E_b'/k)[\psi_b]\longrightarrow0,\\
0\longrightarrow E_b'(k)/\phi_bE_b(k)
 &\longrightarrow\operatorname{Sel}^{\phi_b}(E_b/k)
 \longrightarrow\Sha(E_b/k)[\phi_b]\longrightarrow0.
\end{aligned}
\tag{A.187.10}
\]
Their rightmost groups vanish by (A.187.8).
If \(\xi\in\Sha(E_b/k)[2]\), then
\(\phi_b\xi\in\Sha(E_b'/k)[\psi_b]=0\), and
then \(\xi\in\Sha(E_b/k)[\phi_b]=0\).
The same argument with the dual maps proves
\(\Sha(E_b'/k)[2]=0\).  A nonzero class of
two-power order would have a nonzero multiple of order
two, proving (A.187.1) for both curves.

\emph{The Prym pull--push maps over \(\Q\).}
For each \(D_{A,B}\), first use the unscaled target
\(E_+:Y^2=X^3+2A^3rX\) of (A.185.3), and set
\(\mathcal R=\operatorname{Res}_{L/\Q}E_+\).
The scaling in (A.185.7) identifies its cohomology
with that of the curve just treated over \(L\).
The Shapiro maps of (A.160.16)--(A.160.17), with \(L\)
in place of that statement's field, give
\[
                  \Sha(\mathcal R/\Q)
                     \simeq\Sha(E_+/L).
\tag{A.187.11}
\]
Specifically, on the coinduced module of functions
\(f:G_\Q\to E_+(\overline\Q)\) satisfying
\(f(hg)=hf(g)\) for \(h\in G_L\), evaluation at
one gives the restriction map on cocycles.  For coset
representatives \(r_i\) and
\(r_i g=h_{i,g}r_{i_g}\), its inverse sends a cocycle
\(c\) to \(z(g)(r_i)=c(h_{i,g})\).
Restriction to each local decomposition group partitions
the cosets into the actual places of \(L\) above that
place; it gives the product of their localizations.
Taking the localization kernels proves (A.187.11),
including the archimedean places.  Hence
\(\Sha(\mathcal R/\Q)[2^\infty]=0\).

Let \(f_{\mathcal P}:\mathcal P\to\mathcal R\) be
the isogeny (A.185.6).  Over \(L\) define
\[
 g_{\mathcal P}:E_+\times E_-\longrightarrow\mathcal P_L,
                  \qquad(Q_+,Q_-)\longmapsto
                           q_+^*Q_+ +q_-^*Q_-.
\tag{A.187.12}
\]
Each pullback lands in \(\mathcal P_L\): its image is
connected, and \(\iota\) acts on it by \(-1\), so
it is in the connected anti-invariant abelian subvariety
\(\operatorname{im}(1-\iota)\).  The two pullbacks
are exchanged by \(r\mapsto-r\).  Their sum therefore
descends to the displayed \(\Q\)-homomorphism
\(\mathcal R\to\mathcal P\).  Using the actual
degree-two maps and their involutions gives
\[
\begin{aligned}
 g_{\mathcal P}f_{\mathcal P}
  &=(q_+^*q_{+*}+q_-^*q_{-*})|_{\mathcal P}\\
  &=(2+j_++j_-)|_{\mathcal P}=[2].
\end{aligned}
\tag{A.187.13}
\]
Here \(j_- =\iota j_+\), while \(\iota=-1\)
on \(\mathcal P\), so the two involution terms cancel.
The same identities give
\(f_{\mathcal P}g_{\mathcal P}=[2]\) on
\(\mathcal R\): the diagonal terms are \([2]\),
and a cross term is both invariant and anti-invariant
under the appropriate involution, hence is killed by
two.  A homomorphism from a connected abelian variety
to a finite group scheme is zero, so each cross term
vanishes.  These are exact integral pull--push maps;
the equality has no omitted isogeny multiplier.

For \(\xi\in\Sha(\mathcal P/\Q)[2^\infty]\),
(A.187.11) makes \(f_{\mathcal P}\xi=0\), so
(A.187.13) gives \(2\xi=0\).  The finite group
\(\Sha(\mathcal P/\Q)[2]\) is a quotient of its
finite two-Selmer group.  This proves (A.187.2).
The precise general finiteness input is
\cite[I, Proposition 6.4 and Remark 6.7, pp.75--77]{BIR_MilneADT};
it follows from finite-set cohomological finiteness in
I, Corollary 4.15, with two included in the inverted set.

\emph{The actual isogeny on the genus-twenty-one Jacobian.}
For \(\mathcal B\) in the statement, A.186 gives
vanishing of every genus-two factor's two-primary
Tate--Shafarevich group.  Equation (A.187.1) gives
the same for its three rational elliptic factors,
and (A.187.2) treats its three Prym factors.  Products
commute with the cohomology and localization maps, so
\[
                       2\Sha(\mathcal B/\Q)[2^\infty]=0.
\tag{A.187.14}
\]
There is a concrete inverse-up-to-integer homomorphism
\(G:\mathcal B\to J_8\) determined by
\[
                   G(F(P))=[N]P,\qquad N=\deg F.
\tag{A.187.15}
\]
Indeed the finite characteristic-zero kernel of \(F\)
has order \(N\), so multiplication by \(N\) kills
it.  The indicated map on the finite quotient therefore
descends uniquely as a homomorphism over \(\Q\).
It satisfies \(GF=[N]\); surjectivity of \(F\)
then also gives \(FG=[N]\).  This specifies the
map on cocycles and on every local cohomology group.
For \(\xi\in\Sha(J_8/\Q)[2^\infty]\),
(A.187.14) implies \(2F\xi=0\), hence
\([2N]\xi=0\).  The odd factor of \(N\) acts
invertibly on a two-primary group, yielding the stated
exponent bound \(2^a\).  The resulting group is a
subgroup of the finite \(\Sha(J_8/\Q)[2^a]\),
proving all finiteness assertions in (A.187.3).

\emph{Finite and continuous Kummer maps.}
For any abelian variety here the finite multiplication
sequence gives
\[
0\longrightarrow A(\Q)/2^nA(\Q)
 \xrightarrow{\kappa_{2^n}}\operatorname{Sel}^{2^n}(A/\Q)
 \longrightarrow\Sha(A/\Q)[2^n]\longrightarrow0.
\tag{A.187.16}
\]
Under the transition from level \(n+a\) to \(n\),
the last map is multiplication by \(2^a\).  For
\(A=J_8\) it is zero by (A.187.3), so every image
lies in the displayed finite Mordell--Weil Kummer
image.  Conversely, the class of any rational point
at level \(n\) is the image of the same rational
point at level \(n+a\).  This proves (A.187.4).

For a compatible system of finite Selmer classes,
its \(\Sha\)-images satisfy \(\alpha_n=2^a\alpha_{n+a}=0\).
Each finite class is therefore the Kummer class of a
unique element of \(J_8(\Q)/2^nJ_8(\Q)\), and
compatibility gives one element of the completion.
Conversely every completed point defines such a
compatible system.  The passage to continuous
cohomology is exact, not an additional ramification
assumption.  At each finite level the continuous
cocycles representing a fixed class form a finite
orbit under zero-cochains in the finite coefficient
group.  A compatible system of classes therefore
has compatible representatives by compactness of
these finite orbits: every finite string is obtained
from a representative at its highest level.
Their inverse limit is a continuous \(T_2J_8\)-cocycle.
Likewise, if all finite classes of a continuous
cocycle are zero, the finite nonempty sets of
cobounding zero-cochains have a compatible inverse
limit, which is one integral zero-cochain.  This
proves both surjectivity and injectivity in the
continuous comparison used in (A.187.5).

All finite Kummer sequences are natural in a
homomorphism of abelian varieties and in a field
embedding.  Applying the actual maps \(F,G\),
the Prym pull--push maps, and restriction to
\(G_{\Q_v}\) at every finite level therefore
proves the asserted compatibility for the full
two-adic Selmer groups.  The rank of the completed
Mordell--Weil group is fourteen by A.186.

Finally, compactification followed by the actual
Abel--Jacobi map
\(C_8\to J_8,\ P\mapsto[P-P_0]\), pulls each
multiplication torsor back to the open curve.
For a full Selmer section, its finite classes are
locally the Kummer classes of those local proper
centres.  At a cusp the compactification kills
horizontal inertia and gives the same statement
for its proper boundary point; at a real place
the class is that of its chosen real component.
Thus all the finite classes belong to the Selmer
groups in (A.187.16), and (A.187.5) proves the
last assertion with its stated scope.
\end{proof}
